\documentclass[10pt,a4paper]{amsart}
\usepackage[margin=19mm]{geometry}
\usepackage{amsmath,amssymb,mathtools}
\usepackage[scr=rsfs,cal=euler]{mathalfa}
\usepackage{xfrac}
\usepackage[unicode,hidelinks,bookmarks=false]{hyperref}
\newcommand{\ie}{i.e.\ }
\newcommand{\cf}{cf.\ }
\newcommand{\resp}{respectively\ }
\newcommand{\Subsection}{\subsection}
\providecommand{\nobreakdash}{\nobreak\hbox{-}}
\setlength{\emergencystretch}{2em}
\multlinegap0pt
\usepackage{bm}
\usepackage{amssymb}
\usepackage[shortlabels]{enumitem}
\usepackage{xfrac}
\usepackage[normalem]{ulem}
\usepackage{xcolor}
\usepackage{mathtools}
\usepackage{stackengine}
\usepackage{tikz}
\newtheorem{lemma}{Lemma}
\newtheorem{pf}{Pf}
\newcommand{\Div}{\operatorname{div}}
\newcommand{\Divx}{\operatorname{div}_x}
\newcommand{\Erel}{\mathcal{E}_{\mathrm{rel}}}
\newcommand{\Grad}{\nabla}
\newcommand{\Gradx}{\nabla_x}
\newcommand{\Gradxz}{\nabla_{x,z}}
\newcommand{\Lambdaz}{\Lambda_0}
\newcommand{\Ov}[1]{\overline{#1}}
\newcommand{\T}{\mathbb{T}}
\newcommand{\dt}{{\,\rm d} t }
\newcommand{\dx}{{\,\rm d} {x}}
\newcommand{\dz}{{\rm d} z}
\newcommand{\lr}[1]{\left( #1 \right)}
\newcommand{\vc}[1]{{\bf #1}}
\newcommand{\vm}{\vc{m}}
\newcommand{\vr}{\varrho}
\newcommand{\vu}{\vc{u}}
\newcommand{\vv}{\vc{v}}
\title{Generalized solution and Weak-Strong uniqueness for a barotropic Euler-Riesz system}
\author{Nilasis Chaudhuri}
\date{Source: arXiv:2608.01872v1 (CC BY 4.0)}
\begin{document}
\maketitle

\noindent\emph{Source and licence.} Generalized solution and Weak-Strong uniqueness for a barotropic Euler-Riesz system. Version arXiv:2608.01872v1 (CC BY 4.0), distributed by arXiv under the Creative Commons Attribution 4.0 International licence, http://creativecommons.org/licenses/by/4.0/, as recorded in the arXiv licence metadata for this exact version and shown on the abstract page https://arxiv.org/abs/2608.01872v1. Full licence notice: \texttt{licenses/arxiv-2608.01872-CC-BY-4.0.txt}.\par\smallskip

\noindent\textit{Editorial scope.} Counted unit: the article's lemma lem:generic (source0.tex lines 2047--2065). The article's complete original proof of it is delivered with it (source0.tex lines 2067--2239, one \texttt{proof} environment of 173 lines). No mathematics is written, repaired or re-derived: the statement and the proof are the article's own lines, in the article's own notation, and no retained line is substituted or re-flowed.  Retained uncounted as the source-local context the proof needs: lines 555--561; lines 890--892; lines 992--995; lines 998--1010; lines 1013--1017; lines 1019--1027; lines 1994--1998.  Nothing was omitted from any retained span, so the package carries no omission record.  No retained line contains a citation, so the wrapper carries no bibliography and no bibliography entry.  No retained line contains a figure environment, an image-inclusion command or drawing code, so no image is delivered and the package declares no asset dependency.  Editorial formatting only, in addition: an AMS \texttt{amsart} preamble that restates the article's own declarations and macros used by the retained lines, namely the environments \texttt{lemma}, \texttt{pf} on separate counters, together with the standard packages those lines need.  The retained environments print in this document as Lemma 1, Pf 1 in order of appearance, whereas the article prints the same items with the numbering of its own full document.  The complete native source of the article as distributed in the arXiv e-print for this exact version is bundled unchanged as \texttt{sources/206651/source0.tex}.
	\begin{equation}\label{CS-ext}
		\begin{cases}
			\Div_{x,z}\!\left( z^{1-2s}\,\Grad_{x,z} U \right) = 0 & \text{in } \Omega,\\[2pt]
			-c_{d,s}\,\displaystyle\lim_{z\to 0} z^{1-2s}\,\partial_z U = \vr - \bar\vr & \text{on } \mathbb{T}^d\times\{0\},
		\end{cases}
		\qquad \lim_{z\to\infty} U(x,z) = 0,
	\end{equation}

	\begin{align}\label{Not:CSext}
			\mathfrak E[f]:=U,
	\end{align}

			\begin{equation}\label{eq:w-ERc}
				\int_0^{\tau}\!\!\int_{\mathbb T^d}\big(\vr\,\partial_t\phi+\vm\cdot\Grad_x\phi\big)\,dx\,dt
				+\int_{\mathbb T^d}\vr_0\,\phi(0,\cdot)\,dx- \int_{\mathbb T^d}\vr(\tau)\,\phi(\tau,\cdot)\,dx=0.
			\end{equation}

		\begin{equation}\label{eq:w-ERm}
			\begin{aligned}
				&\int_0^\tau\!\!\int_{\mathbb T^d}\Big(\vm\cdot\partial_t\pmb\psi
				+\big(\tfrac{\vm\otimes\vm}{\vr}+p(\vr)\mathbb I_d\big):\Grad_x\pmb\psi\Big)\,dx\,dt\\
				&\quad -\,c_{d,s}\int_0^\tau\!\!\iint_\Omega \mathbb T_{\rm ER}[U]:\Grad_x\pmb\psi\,dx\,dz\,dt
				+\,\bar\vr\int_0^\tau\!\!\int_{\mathbb T^d}\Lambda_0(U)\,\Div_x\pmb\psi\,dx\,dt\\
				&\quad +\int_0^\tau\!\!\int_{\mathbb T^d}\Grad_x\pmb\psi:d\mathcal R_{fluid}
				-\,c_{d,s}\int_0^\tau\!\!\iint_{\overline\Omega}\Grad_x\pmb\psi:d\mathcal R_{ext}\\
				&\quad +\int_{\mathbb T^d}\vr_0\vu_0\cdot\pmb\psi(0,\cdot)\,dx
				-\int_{\mathbb T^d}\vm(\tau)\cdot\pmb\psi(\tau,\cdot)\,dx
				= 0 .
			\end{aligned}
		\end{equation}

			\begin{equation}\label{eq:w-ext}
				\iint_\Omega z^{1-2s}\Grad_{x,z}U\cdot\Grad_{x,z}\Psi\,dx\,dz
				=c_{d,s}^{-1}\int_{\mathbb T^d}(\vr-\bar\vr)\Lambda_0(\Psi)\,dx
				\qquad\forall\,\Psi\in\dot C_c^\infty(\overline\Omega).
			\end{equation}

			\begin{equation}\label{eq:ei}
				\begin{aligned}
					&\int_{\mathbb T^d}\Big(\tfrac12\tfrac{|\vm|^2}{\vr}+P(\vr)\Big)(\tau)\,dx
					+\tfrac{c_{d,s}}2\iint_\Omega z^{1-2s}|\Grad_{x,z}U(\tau)|^2\,dx\,dz
					+\int_{\mathbb T^d}d\mathcal E_{fluid}(\tau)+\iint_{\overline\Omega}d\mathcal E_{ext}(\tau)\\
					&\le \int_{\mathbb T^d}\Big(\tfrac12\tfrac{|\vm_0|^2}{\vr_0}+P(\vr_0)\Big)\,dx
					+\tfrac{c_{d,s}}2\iint_\Omega z^{1-2s}|\Grad_{x,z}U_0|^2\,dx\,dz .
				\end{aligned}
			\end{equation}

		\begin{equation}\label{eq:RE}
			\Erel(\tau)
			:= \int_{\mathbb T^d}\Big( \frac12\,\vr\,\left| \frac{\vm}{\vr} - \vv\right|^2 + P(\vr\,|\,r) \Big)(\tau)\,dx
			\;+\; \frac{c_{d,s}}{2}\iint_\Omega z^{1-2s}\,\big|\Gradxz (U-V)\big|^2(\tau)\,dx\,dz,
		\end{equation}

	\begin{lemma}[Relative energy inequality I]\label{lem:generic}
		Let $r\in C^1([0,T]\times\mathbb T^d)$ with $r\ge\underline r>0$, $\vv\in C^1([0,T]\times\mathbb T^d;\mathbb R^d)$, $\Ov{r}=\Ov{\vr}$ and let $V=\mathfrak{E}(r-\bar\vr)$ following \eqref{CS-ext} and \eqref{Not:CSext}. Then for a.e.\ $\tau\in(0,T)$ it holds
	\begin{equation}\label{REI1}
		\begin{aligned}
			\Erel(\tau)+&\int_{\mathbb T^d}d\mathcal E_{fluid}(\tau)+\iint_{\overline\Omega}d\mathcal E_{ext}(\tau)
			\\
			\ \le\ \Erel(0)&-\int_0^\tau\!\!\int_{\mathbb T^d}\vr \lr{\tfrac{\vm}{\vr}-\vv}\otimes\lr{\tfrac{\vm}{\vr}-\vv}:\Gradx\vv \dx\dt 
			-\int_0^\tau\!\!\int_{\mathbb T^d}p(\vr\,|\,r)\Divx\vv\dx \dt \\
			&\quad-\int_0^\tau\!\!\int_{\mathbb T^d}\vr\lr{\tfrac{\vm}{\vr}-\vv} \cdot\Big(\partial_t\vv+\vv\cdot\Gradx\vv+\tfrac1r\Gradx p(r)+\Gradx\Lambdaz(V)\Big) \dx \dt \\
			&\quad+\int_0^\tau\!\!\int_{\mathbb T^d}\Big(\tfrac{r-\vr}{r}p'(r)+\Lambdaz(V)-\Lambdaz(U)\Big)\big(\partial_t r+\Divx(r\vv)\big) \dx \dt \\
			&\quad+c_{d,s}\int_0^\tau\!\!\iint_\Omega\T[U]:\Gradx\vv \dx\, \dz\,  \dt
			-\bar\vr\int_0^\tau\!\!\int_{\mathbb T^d}\Lambdaz(U)\Divx\vv \dx \dt \\
			&\quad-\int_0^\tau\!\!\int_{\mathbb T^d}(\vr-r)\,\vv\cdot\Gradx\Lambdaz(V) \dx \dt
			+\int_0^\tau\!\!\int_{\mathbb T^d} \Divx\lr{r\,\vv} \Lambdaz(U) \dx \dt \\
			&\quad-\int_0^\tau\!\!\int_{\mathbb T^d}\Gradx\vv:d\mathcal R_{fluid} \dx \dt 
			+c_{d,s}\int_0^\tau\!\!\iint_{\overline\Omega}\Gradx\vv:d\mathcal R_{ext} \dx \dt = \sum_{i=1}^{11}\mathcal{T}_i
		\end{aligned}
	\end{equation}
	\end{lemma}

	\begin{pf}
		For a.e. $\tau \in (0,T)$, using $p(r)=rP'(r)-P(r)$ we expand \eqref{eq:RE} and it gives
		\begin{align*}
					\Erel(\tau)
				&=\int_{\mathbb T^d}\lr{\!\tfrac12 \frac{|\vm|^2}{\vr} +P(\vr)} \dx 
				-\int_{\mathbb T^d}\!\mathbf m\cdot\vv dx \dt 
				+\int_{\mathbb T^d}\!\vr\,\phi\;  \dx 
				+\int_{\mathbb T^d}\!p(r)\dx  \\
				&	\quad	+\tfrac{c_{d,s}}2\iint_\Omega\! z^{1-2s}|\Gradxz U|^2 \dx \dz 
				-c_{d,s}\!\iint_\Omega\! z^{1-2s}\Gradxz U\cdot\Gradxz V \dx \dz \\
				&\quad
				+\tfrac{c_{d,s}}2\iint_\Omega\! z^{1-2s}|\Gradxz V|^2 \dx \dz=\sum_{i=1}^{7} \mathcal{B}_i (\tau)
		\end{align*}
		with $\phi:=\tfrac12|\vv|^2-P'(r)$, all at time $\tau$.

		
		\smallskip
		\noindent\emph{The terms $\mathcal B_1+\mathcal B_5$.} Their sum is the total energy of the weak solution at time $\tau$, so the energy inequality \eqref{eq:ei} applies verbatim:
		\begin{equation}\label{B15}
			\mathcal B_1(\tau)+\mathcal B_5(\tau)+\int_{\mathbb T^d}d\mathcal E_{fluid}(\tau)+\iint_{\overline\Omega}d\mathcal E_{ext}(\tau)
			\ \le\ \mathcal B_1(0)+\mathcal B_5(0).
		\end{equation}
		This is the only inequality in the computation; all remaining steps are identities.

		\noindent\emph{The term $\mathcal B_2(\tau)$.} We use the momentum identity \eqref{eq:w-ERm} with $\pmb\psi=\vv$, admissible since $\vv\in C^1_{t,x}$.  This yields
		\begin{equation}\label{B2}
			\begin{aligned}
				\big[\mathcal B_2\big]_0^\tau
				&=-\int_0^\tau\!\!\int_{\mathbb T^d}\lr{\vm\cdot\partial_t\vv+\tfrac{\vm\otimes\vm}{\vr}:\Gradx\vv+p(\vr)\Divx\vv}\dx\dt\\
				&\quad+c_{d,s}\int_0^\tau\!\!\iint_\Omega\T_{\rm ER}[U]:\Gradx\vv \dx\dz\dt
				-\bar\vr\int_0^\tau\!\!\int_{\mathbb T^d}\Lambdaz(U)\Divx\vv \dx\dt\\
				&\quad-\int_0^\tau\!\!\int_{\mathbb T^d}\Gradx\vv:d\mathcal R_{fluid} \dx\dt
				+c_{d,s}\int_0^\tau\!\!\iint_{\overline\Omega}\Gradx\vv:d\mathcal R_{ext} \dx\dt .
			\end{aligned}
		\end{equation}

		\noindent\emph{The terms $\mathcal B_3+\mathcal B_4+\mathcal B_4$.} The weak continuity equation \eqref{eq:w-ERc} with $\phi=\tfrac12|\vv|^2-P'(r)\in C^1_{t,x}$, together with $\displaystyle \left[\int_{\T^d} p(r) \dx  \right]_{t=0}^{t=\tau}=\int_0^\tau\!\int_{\T^d}\partial_t p(r) \dx \dt$, gives
		\begin{equation}\label{B34}
			\begin{aligned}
			\big[\mathcal B_3+\mathcal B_4\big]_0^\tau
			=&\int_0^\tau\!\!\int_{\mathbb T^d}\lr{\vr\,\partial_t\tfrac12|\vv|^2+\vm\cdot\Gradx\tfrac12|\vv|^2}\dx\dt \\
			&+\int_0^\tau\!\!\int_{\mathbb T^d}\lr{-\vr\,\partial_tP'(r)-\vm\cdot\Gradx P'(r)+\partial_tp(r)}\dx\dt .
		\end{aligned}
		\end{equation}
		
	
		Next, We put a bit more focus on the Terms $\mathcal{B}_6$ and $\mathcal{B}_7$

		
		Here we use that $V=\mathfrak{E}(r-\bar\vr)$ with $\bar\vr=\bar r$. Note that, no equation for $(r,\vv)$ is invoked at this stage. 
		

\emph{The term $\mathcal B_6(\tau)$.} Applying \eqref{eq:w-ext} with $\Psi=V$, and using
$\int_{\mathbb T^d}\Lambdaz(V)\,\dx=0$,
\begin{equation*}
	\mathcal B_6(\tau) = -\int_{\mathbb T^d}(\vr-\bar\vr)\,\Lambdaz(V)\,\dx
	= -\int_{\mathbb T^d}\vr\,\Lambdaz(V)\,\dx .
\end{equation*}
Since $r\in C^1((0,T)\times\mathbb T^d)$, the trace $\Lambdaz(V)\in C^1((0,T)\times\mathbb T^d)$.
Testing the weak continuity equation \eqref{eq:w-ERc} with $\phi=\Lambdaz(V)$ gives
\begin{equation}\label{B6-dt}
	\mathcal B_6(\tau)-\mathcal B_6(0)
	= -\int_0^\tau\!\!\int_{\mathbb T^d}\Big(\vr\,\partial_t\Lambdaz(V)
	+ \vm\cdot\Gradx\Lambdaz(V)\Big)\,\dx\,\dt .
\end{equation}

\smallskip
\noindent\emph{The term $\mathcal B_7(\tau)$.} Testing \eqref{eq:w-ext} with $\Psi=V$ once more
--- now reading it as the energy identity --- yields
\begin{equation}\label{B7-static}
	\mathcal B_7 = \tfrac{c_{d,s}}2\iint_\Omega z^{1-2s}\,|\Gradxz V|^2\,\dx\,\dz
	= \tfrac12\int_{\mathbb T^d}(r-\bar\vr)\,\Lambdaz(V)\,\dx .
\end{equation}
Since $E$ is linear and $t$-independent, $\partial_t V=E(\partial_t r)$, and $\partial_t r$ is
mean-zero as $\bar\vr$ is constant; differentiating \eqref{B7-static} in time and testing
\eqref{eq:w-ext} with $\Psi=\partial_t V$,
\begin{equation}\label{B7-symm}
	\frac{d}{dt}\,\mathcal B_7(t)
	= c_{d,s}\iint_\Omega z^{1-2s}\,\Gradxz V\cdot\Gradxz\partial_t V\,\dx\,\dz
	= \int_{\mathbb T^d}\partial_t r\,\Lambdaz(V)\,\dx ,
\end{equation}
the factor $\tfrac12$ absorbed because the two symmetric contributions of the product rule
coincide. Hence
\begin{equation}\label{B7-dt}
	\mathcal B_7(\tau)-\mathcal B_7(0)
	= \int_0^\tau\!\!\int_{\mathbb T^d}\partial_t r\,\Lambdaz(V)\,\dx\,\dt .
\end{equation}

It is convenient to combine $\mathcal B_6$ and $\mathcal B_7$. Since $\mathfrak{E}$ is linear and
$t$-independent, $\partial_t V=\mathfrak{E}(\partial_t r)$ is the extension of the mean-zero source
$\partial_t r$. Testing \eqref{eq:w-ext} for $U$ (source $\vr-\bar\vr$) with $\Psi=\partial_t V$,
and \eqref{eq:w-ext} for $\partial_t V$ (source $\partial_t r$) with $\Psi=U$, both have the same
left-hand side $\iint_\Omega z^{1-2s}\Gradxz U\cdot\Gradxz\partial_t V\,\dx\dz$; equating the
right-hand sides,
\[
\int_{\mathbb T^d}(\vr-\bar\vr)\,\partial_t\Lambdaz(V)\,\dx
= \int_{\mathbb T^d}\partial_t r\,\Lambdaz(U)\,\dx ,
\]
and since $\int_{\mathbb T^d}\partial_t\Lambdaz(V)\,\dx=0$, also
$\int_{\mathbb T^d}\vr\,\partial_t\Lambdaz(V)\,\dx=\int_{\mathbb T^d}\partial_t r\,\Lambdaz(U)\,\dx$.
Adding \eqref{B6-dt} and \eqref{B7-dt}, the two time derivatives of $\Lambdaz(V)$ cancel and there
remains
\begin{equation}\label{B67}
	\big[\mathcal B_6(s)+\mathcal B_7(s)\big]_{s=0}^{s=\tau}
	= -\int_0^\tau\!\!\int_{\mathbb T^d}\vm\cdot\Gradx\Lambdaz(V)\,\dx\,\dt
	+ \int_0^\tau\!\!\int_{\mathbb T^d}\partial_t r\,\big(\Lambdaz(V)-\Lambdaz(U)\big)\,\dx\,\dt .
\end{equation}


		We rewrite the second term in the R.H.S of \eqref{B67} as 

		\begin{align*}
			&	\int_{\mathbb T^d}\partial_t r\,\big(\Lambdaz(V)-\Lambdaz(U)\big) \dx \\
			&= \int_{\mathbb T^d}\big(\Lambdaz(V)-\Lambdaz(U)\big)\big(\partial_t r+\Divx(r\vv)\big) \dx
			+ \int_{\mathbb T^d}\big( r\,\vv\cdot \Gradx \Lambdaz(V)+ \Div \lr{ r\,\vv}\Lambdaz(U)\big) \dx .
		\end{align*}
The two terms are treated differently on purpose: $\Lambdaz(V)\in C^1_{t,x}$ carries a derivative, whereas $\Lambdaz(U)\in\dot H^{s}(\mathbb T^d)$ does not, and the pairing is kept in the form in which it is a Lebesgue integral.

		
	\smallskip
	
		\noindent\emph{Regrouping.} All identities below are pointwise a.e., with the convention that expressions carrying $\vr$ in the denominator vanish on $\{\vr=0\}$, where also $\vm=0$; note $\vm-\vr\vv=0$ there as well, so both sides vanish on the vacuum set. %Recall the index convention $(\Gradx\vv)_{jk}=\partial_{x_j}v_k$, $A:B=A_{jk}B_{jk}$.
	
	The terms of \eqref{B2}, \eqref{B34} carrying derivatives of $\vv$ combine as follows. From $\vr\,\partial_t\tfrac12|\vv|^2=\vr\,\vv\cdot\partial_t\vv$,
	\[
	-\vm\cdot\partial_t\vv+\vr\,\partial_t\tfrac12|\vv|^2=-(\vm-\vr\vv)\cdot\partial_t\vv ,
	\]
	while $\vm\cdot\Gradx\tfrac12|\vv|^2=(\vm\otimes\vv):\Gradx\vv$ gives
	\begin{align*}
		-\tfrac{\vm\otimes\vm}{\vr}:\Gradx\vv+\vm\cdot\Gradx\tfrac12|\vv|^2
		&=-\tfrac{\vm\otimes(\vm-\vr\vv)}{\vr}:\Gradx\vv\\
		&=-\tfrac{(\vm-\vr\vv)\otimes(\vm-\vr\vv)}{\vr}:\Gradx\vv-(\vm-\vr\vv)\cdot\lr{\vv\cdot\Gradx\vv},
	\end{align*}
	where the last step used $\vm=(\vm-\vr\vv)+\vr\vv$ and $\big(\vv\otimes(\vm-\vr\vv)\big):\Gradx\vv=(\vm-\vr\vv)\cdot(\vv\cdot\Gradx\vv)$. Hence
	\begin{equation}\label{grp-kin}
		\begin{aligned}
			&	-\vm\cdot\partial_t\vv-\tfrac{\vm\otimes\vm}{\vr}:\Gradx\vv
			+\vr\,\partial_t\tfrac12|\vv|^2+\vm\cdot\Gradx\tfrac12|\vv|^2 \\
			&	=-\tfrac{(\vm-\vr\vv)\otimes(\vm-\vr\vv)}{\vr}:\Gradx\vv
			-(\vm-\vr\vv)\cdot\lr{\partial_t\vv+\vv\cdot\Gradx\vv} .			
		\end{aligned}
	\end{equation}
	The first term in the R.H.S of \eqref{grp-kin} is exactly leads to the term $\mathcal{T}_2$ in \eqref{REI1}.
	For the pressure terms we use $rP''(r)=p'(r)$, i.e.\ $\partial_tP'(r)=\tfrac{p'(r)}{r}\partial_tr$ and $\Gradx P'(r)=\tfrac1r\Gradx p(r)$, and split $\vm=(\vm-\vr\vv)+\vr\vv$ in the flux term. Then
	\begin{equation}\label{grp-pres}
		\begin{aligned}
			&\int_{\mathbb T^d}\lr{-p(\vr)\Divx\vv-\vr\,\partial_tP'(r)-\vm\cdot\Gradx P'(r)+\partial_tp(r)}\dx\\
			&\qquad=\int_{\mathbb T^d}\lr{-p(\vr\,|\,r)\Divx\vv-(\vm-\vr\vv)\cdot\tfrac1r\Gradx p(r)
				+\tfrac{r-\vr}{r}p'(r)\big(\partial_tr+\Divx(r\vv)\big)}\dx .
		\end{aligned}
	\end{equation}
	Clearly the first term in the  R.H.S of \eqref{grp-pres} is $\mathcal{T}_3$ in \eqref{REI1}.

	
	\smallskip


	 \noindent\emph{Conclusion.} In \eqref{B67} we split $\vm=(\vm-\vr\vv)+\vr\vv$ in the first term. Together with the identity above for $\partial_t r$, the contribution of $\mathcal B_6+\mathcal B_7$ thus consists of $-(\vm-\vr\vv)\cdot\Gradx\Lambdaz(V)$, of the continuity residual with coefficient $\Lambdaz(V)-\Lambdaz(U)$, and of the remainder
	\begin{align*}
			 -\int_{\mathbb T^d}\vr\,\vv\cdot\Gradx\Lambdaz(V)\dx 
		 + \int_{\mathbb T^d}\big( r\,\vv\cdot \Gradx \Lambdaz(V)+ \Div \lr{ r\,\vv}\Lambdaz(U)\big) \dx\\
		=-\int_{\mathbb T^d}(\vr-r)\,\vv\cdot\Gradx\Lambdaz(V)\dx+\int_{\mathbb T^d} \Div\lr{r\,\vv} \Lambdaz(U)\dx ,
	\end{align*}
	
	 which is the equals to $\mathcal{T}_8+\mathcal{T}_9$ in \eqref{REI1}. \par 
	 Next, adding now \eqref{B15}, \eqref{B2}, \eqref{B34} in the form \eqref{B67}, \eqref{grp-kin} and \eqref{grp-pres}, the three terms proportional to $\vm-\vr\vv$ assemble the momentum residual $\mathcal{T}_4$ in \eqref{REI1}, i.e. 
	 \[
	 -\int_0^\tau\!\!\int_{\mathbb T^d}(\vm-\vr\vv)\cdot\Big(\partial_t\vv+\vv\cdot\Gradx\vv+\tfrac1r\Gradx p(r)+\Gradx\Lambdaz(V)\Big)\dx\dt .
	 \]
	 Meanwhile the residual terms of \eqref{grp-pres} and of \eqref{B67} assemble the continuity residual with coefficient $\tfrac{r-\vr}{r}p'(r)+\Lambdaz(V)-\Lambdaz(U)$, that is the term $\mathcal{T}_5$ in \eqref{REI1} . Hence, collecting all the information from above we obtain the relative energy inequality \eqref{REI1}.


	\end{pf}

\end{document}
